<<<P001027>%
Рассмотрим <<<P000915> <<<p000010> в форме

<<P001185>
<<P000933>
соответственно

<<P001186>
<<P000934>.

\noindent{%}
<<<P000015> быть суперклассом замкнутого класса <<<P000016>.
Ошибка времени компиляции в классе <<<P000017> не объявлять
Генеративный конструктор по имени <<<P000018> (перенаправлено с «P000935»).
В противном случае, статический анализ <<<P000020> выполняется
как указано в разделе ~<<<P000806>,
Как будто <<P000936> соответственно <<P000937>
имел функциональный тип обозначенного конструктора,
% % TODO(Eernst): Следующее очень неточно, оно просто служит для запоминания.
% %, что мы должны указать, как обращаться с переменными типа в этом параметре
% % часть. Одна вещь, которую мы лелеем, заключается в том, что суперкласс может быть смесью.
%-ное применение.
Замена формальных переменных типа суперкласса
для соответствующего фактического типа аргументов, переданных суперклассу
в заголовке текущего класса.

<<<P001028>%
Давай. <<<P001188>{k} — генеративный конструктор.
<<<P000021> может включать в свой список инициализаторов не более одного суперинициализатора
или возникает ошибка времени компиляции.
Если нет суперинициализатора,
добавлен неявный суперинициализатор формы \SUPER{}()
в конце списка инициализаторов <<<P000022>,
Если только класс не является классом <<P000938>.
Это ошибка времени компиляции, если появляется суперинициализатор.
в Список инициализаторов P000023 в любой другой позиции, чем в конце.
Ошибка времени компиляции, если более одного инициализатора соответствует
Для данного экземпляра переменная появляется в списке инициализаторов <<<P000024>.
Это ошибка времени компиляции, если список инициализаторов <<<P000025> содержит
инициализатор переменной, инициализированный посредством
форма инициализации <<<P000026>.
Это ошибка времени компиляции, если список инициализаторов <<<P000027> содержит
инициализатор для конечной переменной <<<P000028> чье заявление включает
Выражение инициализации.
Ошибка времени компиляции <<<P000029> включает в себя инициализацию формального
для конечной переменной <<<P000030> чье заявление включает
Выражение инициализации.

<<<P001029>%
Пусть <<<P001190>{f} будет переменной конечного экземпляра, объявленной в
Немедленно включающий класс или число.
Ошибка времени компиляции происходит, если <<<P000031> инициализируется
одним из следующих средств:

<<P000770>
<<P001191> <<<P000032> объявляется инициализацией переменной декларации.
<<P001192> <<P000033> инициализируется посредством инициализации формы <<<P000034>.
<<P001193> <<P000035> имеет инициализатор в списке инициализаторов <<<P000036>.
<<P000788>

<<<P001030>%
Это ошибка времени компиляции, если список инициализаторов <<<P000037> содержит
инициализатор для переменной, которая не
переменная экземпляра, объявленная в непосредственно окружающем классе.

<<<P001194>{%>
Список инициализаторов может содержать инициализатор для
любой переменной экземпляра, объявленной непосредственно окружающим классом,
Даже если это не окончательно.%
?

<<<P001031>%
Это ошибка времени компиляции, если генеративный конструктор класса <<<P000939>
Включает в себя суперинициализатор.


<<<P001195> Выполнение генеративных конструкторов.
<<P001154>

<<<P001032>%
<<<P001174>%
Исполнение генеративного конструктора <<P000038> Тип <<P000039>
для инициализации нового экземпляра <<<P000040>
всегда выполняется в отношении набора привязок для его формальных параметров.
и типовые параметры непосредственного замкнутого класса или перечня, связанного с
набор фактических аргументов типа <<<P000041>, \DefineSymbol{\List{t}{1}{m}}.

<<<P001198>{%>
Эти связывания обычно определяются выражением создания экземпляра.
Строитель (прямо или косвенно).
Тем не менее, они также могут быть определены отражающим вызовом.
?

<<<P001033>%
Если <<P000042> перенаправляет, тогда его пункт перенаправления имеет форму

<<P001199>
<<P000940>
где \DefineSymbol{g} идентифицирует другой генеративный конструктор
непосредственно окружающего класса.
Затем выполняется <<<P000047> Для инициализации <<<P000048> доходы от
Оценка списка аргументов
<<P000941>
Список аргументов <<<P000052> в форме
<<P000942>
В среде, где
Типовые параметры ограждающего класса связаны с
<<P000056>.

<<<P001034>%
Далее, тело <<<P000057> Выполняется для инициализации <<<P000058>
В отношении переплетов, которые составляют карту
Формальные параметры <<<P000059> для соответствующих объектов
в реальном списке аргументов <<<P000060>,
<<<<P001201>< Ссылаясь на <<P000061>,
и типовые параметры непосредственного замкнутого класса или перечня, связанного с
\List{t}{1}{m}.

<<<P001035>%
В противном случае <<<P000062> не перенаправляется.
После этого казнь осуществляется следующим образом:

<<<P001036>%
Заявления о переменных экземплярах непосредственно прилагаемого класса или перечня
Посещаются в том порядке, в котором они появляются в тексте программы.
Для каждого такого заявления <<<P000063>, если <<<P000064> имеет форму
\code{\synt{финалConstVarOrType) <<<P000065> = <<P000066>;
<<<P000067> оценивается на предмет <<P000068>
и переменная экземпляра $v$ <<<P000070> Ссылка: <<<P000071>.

<<<P001037>%
Любые инициализирующие формы, объявленные в списке параметров <<<P000072>
выполняются в том порядке, в котором они появляются в тексте программы.
% На самом деле этот порядок ненаблюдаем, это можно сделать в любое время.
% до начала работы организма, так как эти <<<P001205>.
Затем инициализаторы списка инициализаторов <<<P000073> выполняются для инициализации <<<P000074>
в порядке их появления в программе, как описано ниже
(p.\,\pageref{executionOfInitializerLists}).

<<<<P001207>{%>
Мы могли бы наблюдать порядок, выполняя побочные действия, называемые внешними процедурами.
Поэтому необходимо указать порядок.%
?

<<<P001038>%
Если есть переменная <<<P000075> объявленный
непосредственно включающий класс или
Он не привязан к объекту,
Все такие переменные инициализируются нулевым объектом (<<P000807>).

<<<P001039>%
Затем, если прилагаемый класс не является <<<P000943>, явно указанный или
Неявно добавленный суперинициализатор (<<<P000808>) выполняется для
дальнейшей инициализации <<P000076>.

<<<P001040>%
После завершения суперинициализатора тело <<<P000077> исполняется
в сфере, где <<<P001208>{} связано с <<<P000078>.

<<<<P001209>{%>
Этот процесс гарантирует, что ни одна неинициализированная переменная конечного экземпляра не
Его всегда можно увидеть по коду.
Обратите внимание, что \THIS{} не находится в области действия на правой стороне инициализатора.
(см. <<P000809>)
Таким образом, ни один метод экземпляра не может выполняться во время инициализации:
метод не может быть использован непосредственно,
не может быть \THIS{} передаваться в любой другой код, на который ссылаются
в инициализаторе.%
?


<<<p001212>> Исполнение списков инициализаторов.
<<P001155>

<<<P001041>%
Во время выполнения генеративного конструктора инициализировать экземпляр
\DefineSymbol{i)
выполнение инициализатора формы <<P000944>
идет следующим образом:

<<<P001042>%
Во-первых, выражение <<P000081> оценивается на предмет <<P000082>.
Затем переменная экземпляра <<<P000083> <<<P000084> Обозначается как <<P000085>.
% Эта ошибка может произойти из-за неявного слепка.
Ошибка динамического типа, если динамический тип <<<P000086> не является
Подтип фактического типа
(<<P000810>)
переменной <<<P000087>.

<<<P001043>%
Исполнение инициализатора, который является утверждением, происходит путем
Выполнение утверждения (<<P000811>).

<<<P001044>%
Рассмотрим суперинициализатор <<<P001214>{s} формы

<<<P001215>
<<P000945>
соответственно

<<<P001216>
<<P000946>.

<<<P001045>%
<<<P001175>%
<<<P000093> класс, в котором <<<P000094> Появляется и позволяет <<<P000095> Это суперкласс P000096.Если <<P000097> является общим (<<<P000812>),
давать <<<P000098> - фактический тип аргументов, переданных <<<P000099>,
Получено путем замены действительных привязок \List{t}{1}{m)
Формальные параметры типа <<<P000100>
в суперклассе, как указано в заголовке <<<P000101>.
Давай. \DefineSymbol{k} быть конструктором заявлено в $S$ названный
<<P000103> соответственно <<P000947>.

<<<P001046>%
Выполнение <<<P000105> идет следующим образом:
Список аргументов
<<P000948>
оценивается по фактическому списку аргументов <<<P000109> в форме
<<P000949>.
Затем корпус суперконструктора <<<P000113> исполняется
в среде, где формальные параметры <<<P000114> связаны с
соответствующие фактические аргументы из <<<P000115>,
Формальные параметры типа <<<P000116> Обозначается как <<P000117>.


<<<P001219> Фабрики
<<P001156>

<<<P001047>%
A <<<P001220>{завод}{конструктор!завод}
Конструктор, предваряемый встроенным идентификатором
(<<P000813>)
<<<P001221>.

<<P000771>
<factoryConstructorSignature> ::= \gnewline{}
<<<P001223>> \FACTORY{{}<constructorName> <formalParameterList>
<<P000789>

<<<P001048>%
Тип возврата фабрики, подпись которой
форма \FACTORY{} <<P000118> или
форма \FACTORY{} <<P000950>
<<<P000120>, если <<<P000121> не является общим типом;
В противном случае тип возврата
<<P000951>
где <<P000124> являются типовыми параметрами ограждающего класса.

<<<P001049>%
Ошибка времени компиляции <<<P000125> Это не имя
% % TODO (Eernst): Come Dart 3.0, add 'mixin'.
Немедленно включающий класс или число.

<<<P001050>%
% Эта ошибка может произойти из-за неявного слепка.
Это динамическая ошибка типа, если завод возвращает ненулевой объект.
тип которого не является подтипом его фактического
(<<P000814>)
Тип возврата.

<<<P001227>{%>
Бесполезно позволять фабрике возвращать нулевой объект (P000815).
Но это более единообразно, чтобы позволить это, как это делают правила в настоящее время.
?

<<<P001228>{%>
Заводы устраняют классические недостатки, связанные с конструкторами
на других языках.
Заводы могут выпускать экземпляры, которые не были недавно выделены:
Они могут прийти из кэша.
Кроме того, фабрики могут возвращать экземпляры разных классов.
?


<<<P001229> Перенаправление заводских строителей.
<<P001157>

<<<P001051>%
A <<<P001230>{перенаправление заводского конструктора}{конструктор!перенаправление завода}
указывает вызов конструктору другого класса, который должен быть использован
Всякий раз, когда вызывается перенаправляющий конструктор.

<<P000772>
<redirectoryConstructorSignature> ::= \gnewline{}
<<<P001232>> \FACTORY{} <constructorName> <formalParameterList> \gnewline{}
<конструкторское проектирование>

<constructorDesignation> ::= <typeIdentifier>
<<P001235> <qualifiedName>
<<<P001236> <typeName> <typeArguments> ('.' <identifier>)?
<<P000790>

Предположим, что
<<<P001176>%
<<P000952>
- наименование и формальные параметры типа ограждающего класса,
<<<P001177>%
<<<P000128> <<<P001237>> или пустой,
<<<P000129> - <<<P000130> или <<<P000131> для некоторого идентификатора <<<P000132>,
и \DefineSymbol{<<<P001239>> является идентификатором,
Затем рассмотрите декларацию перенаправляющего заводского конструктора.
<<<P001240>{k} одной из форм

<<<P000001>

<<P001241>
<<<P001178>%
где <<P000133> является одной из форм
<<P000953> или
<<P000954>.

<<<P001052>%
Ошибка времени компиляции <<<P000138> не обозначает
класс, доступный в текущем объеме.
Если <<P000139> обозначает такой класс <<<P000140>,
Ошибка времени компиляции <<<P000141> не обозначает конструктора.
% Этой индукции достаточно, чтобы проверить абстрактность на один уровень ниже.
%, потому что это ошибка перенаправления, если это происходит после нескольких
% перенаправления:
В противном случае это ошибка времени компиляции.Если <<P000142> обозначает генеративный конструктор и <<<P000143> является абстрактным.
В противном случае,
<<<P001242>{redirectee constructor}{constructor!redirectee}}
Для этого декларация является конструктором <<<P001243>{k'} обозначается <<<P000144>.

<<<P001053>%
Перенаправляющий заводской конструктор <<<P000145> <<<P000916>
от перенаправляющего заводского конструктора <<P000146> если {f}
<<P000147> является конструктором перенаправления <<<P000148>,
<<<P000149> Конструктор перенаправления <<<P000150>
<<<P000151> Перенаправление доступно из <<<P000152>.
Ошибка времени компиляции при перенаправлении конструктора завода
Это перенаправление-доступно от самого себя.

<<<P001054>%
Пусть <<<P000153> Статический тип списка аргументов
(<<P000816>)
<<P000955>
когда <<P000155> не приводит никаких аргументов, и
<<P000956>
когда <<P000157> Приведем несколько названных аргументов.
Ошибка времени компиляции <<<P000158>
не является подтипом для формального списка параметров перенаправления.

<<<P001244>{%>
Требуется подтип матча
(в отличие от более прощального матча)
который используется с конструктором генеративного перенаправления,
потому что завод перенаправляет конструктора <<<P000159> всегда призывает
его перенаправление <<P000160> с
Те же самые аргументы, что и <<P000161> получено.
Это означает, что понижение фактического аргумента
<<<P000162> и <<P000163>
не используется, так как фактический аргумент
тип, требуемый <<<P000164>,
или это будет означать динамическую ошибку, которая просто задерживается на один шаг.
?

<<<P001245>{%>
Обратите внимание, что необычный случай охватывается
<<<P000165> или <<P000166> или оба равны нулю,
в этом случае формальный список параметров типа класса <<<P000167>
и/или фактический список аргументов типа
Конструктор перенаправления опущен
(<<P000817>)%
?

<<<P001055>%
Ошибка времени компиляции <<<P000168> явно указывает
значение по умолчанию для необязательного параметра.

<<<P001246>{%>
Значения по умолчанию, указанные в <<<P000169> будет игнорироваться,
Так как это параметры <<<P001247><<<P000170>>.
Следовательно, значения по умолчанию не допускаются.%
?

<<<P001056>%
Ошибка времени компиляции, если формальный параметр <<P000171> имеет значение по умолчанию
тип которого не является подтипом аннотации типа
по соответствующему формальному параметру в <<<P000172>.

<<<P001248>{%>
Обратите внимание, что невозможно изменить аргументы, передаваемые <<<P000173>.%
?

<<<P001249>{%>
На первый взгляд, можно подумать, что
Обычные заводские конструкторы могли бы просто создать
других классов и вернуть их,
Перенаправлять фабрики не нужно.
Однако перенаправление заводов имеет ряд преимуществ:

<<P000773>
<<P001250>
Абстрактный класс может быть постоянным конструктором.
Использует константный конструктор другого класса.
<<P001251>
Перенаправляющий завод-конструктор избегает необходимости в экспедиторах
повторять формальные параметры и их значения по умолчанию.
<<<P000791>%
?

<<<P001057>%
Ошибка времени компиляции <<<P000174> Приставка с приставкой <<<P001252>{} модификатор
<<<P000175> Не является конструктором (P000818).

<<<P001058>%
Пусть \DefineSymbol{\List T}{1}{m}} — фактический тип аргументов, переданный <<<P000176>
в декларации <<<P000177>.
Пусть \DefineSymbol{<<P001256>>>{X}{1}{m}} являются формальными параметрами типа, объявленными
Класс, содержащий декларацию <<<P000178>.
<<<P001257> F' - тип функции <<<P000179> (P000819).
Ошибка времени компиляции <<<P000180>
не является подтипом функции типа <<<P000181>.

<<<P001258>{%>
В случае, когда два класса не являются общими
Это просто проверка подтипа на типы функций двух конструкторов.
В общем, это означает, что полученный объект соответствует% % TODO: Come Dart 3.0, добавьте «миксин».
Интерфейс класса или перечня, который непосредственно включает <<<P000182>.%
?

<<<P001059>%
Для динамической семантики,
предположить, что <<<P001259>{k} - перенаправляющий заводской конструктор
и \DefineSymbol{k'} является перенаправлением <<<P000183>.

<<<P001060>%
% Эта ошибка может произойти из-за неявного слепка.
Это ошибка динамического типа, если фактический аргумент принят в вызове <<<P000184>
не является подтипом фактического типа (<<P000820>)
соответствующего формального параметра в декларации <<<P000185>.

<<<P001061>%
При перенаправлении <<<P000186> является фабрикантом,
Выполнение <<<P000187> Выполнение <<<P000188>
Фактические аргументы перешли к <<<P000189>.
Результат исполнения <<<P000190> является результатом <<<P000191>.

<<<P001062>%
При перенаправлении <<<P000192> является генеративным конструктором,
<<<P000193> Будьте свежим примером (<<P000821>)
Класс, содержащий <<<P000194>.
Выполнение <<<P000195> Тогда это равносильно исполнению <<<P000196> инициализировать <<P000197>,
регулируется теми же правилами, что и выражение создания экземпляра
(<<P000822>).
Если <<P000198> Затем выполняется обычное выполнение <<<P000199>
обычно возвращается <<P000200>,
иначе <<P000201> Завершает, бросая исключение и след стека
Выброшено <<<P000202>.


<<<P001261> Постоянные строители.
<<P001158>

<<<P001063>%
A <<<P001262>{постоянно} Конструктор (Constructor!Constant)
Может использоваться для создания объектов постоянной времени компиляции (<<P000823>).
Константный конструктор закрепляется зарезервированным словом <<<P001263>.

<<P000774>
<constantConstructorSignature> ::= \gnewline{}
\CONST{{}<constructorName> <formalParameterList>
<<P000792>

<<<P001266>{%>
Постоянные конструкторы имеют более сильные ограничения, чем другие.
Например, вся работа
Конструктор генеративной постоянной без перенаправления
должно быть выполнено в его инициализаторах
При инициализации выражений
Примерные переменные закрывающего класса
(Возможно, это произошло раньше).
Потому что инициализация должна быть постоянной.%
?

<<<P001064>%
Постоянные перенаправления генеративных и заводских конструкторов указаны в другом месте.
(p.\,\pageref{redirectingGenerativeConstructors}}
p.\,\pageref{redirectingFactoryConstructors}.
Отныне этот раздел посвящен
Неперенаправляющие константные генераторы.

<<<P001065>%
Это ошибка времени компиляции, если конструатор генеративной константы не перенаправляет
Объявляется классом, который имеет переменную экземпляра, которая не является окончательной.

<<<P001269>{%>
Вышеуказанное относится как к локально объявленным, так и к унаследованным переменным экземпляров.
?

<<<P001066>%
Если неперенаправленный константный конструктор \DefineSymbol{k)
Объявляется классом <<<P000203>,
Это ошибка времени компиляции
переменная, объявленная в <<<P000204>
иметь инициализирующее выражение, которое не является постоянным выражением.

<<<P001271>{%>
Суперкласс <<<P000205> не может иметь такого инициализирующего выражения <<<P000206> тоже.
Если у него есть неперенаправляющий константный конструктор
<<<P000207> Это ошибка,
Если у вас нет такого конструктора
затем (явный или неявный) суперинициализатор в <<<P000208> Ошибка.%
?

<<<P001067>%
Суперинициализатор, который появляется, явно или неявно,
в списке инициализаторов постоянного конструктора
Необходимо указать конструктор генеративной постоянной
Сверхкласс немедленно закрывающегося класса,
или возникает ошибка времени компиляции.

<<<P001068>%
Любое выражение, которое появляется внутри
Список инициализаторов постоянного конструктора
Должно быть потенциально постоянное выражение
(<<P000824>),
или возникает ошибка времени компиляции.

<<<P001069>%
Когда константный конструктор \DefineSymbol{k) На него ссылаются
Постоянное выражение объекта,Ошибка времени компиляции, если
Обращение к <<P000209> Во время пробежки было бы исключение,
Ошибка времени компиляции, если
Подмена фактических аргументов формальными параметрами
дает инициализирующее выражение <<<P000210> в списке инициализаторов <<<P000211>
Это не постоянное выражение.

<<<P001273>{%>
Например, если <<<P000212> <<<P000957>
где <<P000958> является формальным аргументом <<<P000213> с типом <<<P001274>,
<<<P000214> потенциально постоянен и может быть использован в списке инициализаторов <<<P000215>.
Ошибка вызова <<<P000216> С аргументом типа <<P000959>
Если <<P000960> Отличается от класса <<<P000961>,
Даже если <<P000962> <<<P000963> Геттер,
и то же выражение будет оцениваться без ошибок во время выполнения.
?


<<<P001275> Статические методы
<<P001159>

<<<P001070>%
<<<P001276>{Статические методы}{метод!статика}
функции, кроме геттеров или сеттеров,
декларация которых непосредственно содержится в классовой декларации
Об этом говорится в статье <<<P001277>.
Статические методы класса <<<P000217> Статические методы, объявленные <<<P000218>.

<<<P001278>{%>
Наследование статических методов мало полезно в Дарте.
Статические методы не могут быть отменены.
Любой необходимый статический метод можно получить из его декларирующей библиотеки.
И нет необходимости вносить его в сферу действия через наследование.
Опыт показывает, что разработчики запутались в
Идея наследственных методов, которые не являются методами экземпляра.

Конечно, вся концепция статических методов является спорной.
Но он сохранился здесь, потому что многие программисты знакомы с ним.
Статические методы Дарта можно рассматривать как функции библиотеки.
?

<<<P001279>{%>
Заявления статического метода могут противоречить другим заявлениям
(<<P000825>)%
?


<<<p001280>> Суперклассы.
<<P001160>

% % TODO(Eernst): Надо сказать, что сверхкласс, который получается
% % по применению смешивания является общим, когда <<<P000219> является общим или, по крайней мере,
%, когда классы используют одну или несколько переменных типа <<<P000220>
% % в \EXTENDS{} или \WITH{} Статья <<P000221>. Ниже говорится, что
% % этих оговорок находятся в области параметров типа <<<P000222>, но это
%% не позволяют нам говорить о суперклассе как о реальном, автономном
Класс %% (если только мы не начнем определять вложенные классы так, что
Процентный суперкласс может быть объявлен в этой области.

<<<P001071>%
Суперкласс <<<P000223> класса <<<P000224> чье заявление имеет пункт
\code{<<P001284>>>{} <<<P000225>>
и пункт о продлении
\code{\EXTENDS{} <<<P000226>>
абстрактный класс, полученный путем применения
Композиция микширования (<<P000826>>>) <<<P000227><<<P000228>.
Имя <<<P000229> является новым идентификатором.
Если нет <<<P001287> Затем указывается пункт <<<P001288>{} пункт
Класс <<<P000230> Определяет свой суперкласс.
Если нет <<<P001289> пункт уточняется, затем либо:

<<P000775>
<<P001290> <<P000231> – это <<P000964>, не имеющий суперкласса. или
<<P001291> Класс <<<P000232> Считается, что имеет <<<P001292> пункт о форме
\code{<<P001294>>>{) Объект, и вышеприведенные правила применяются.
<<P000793>

<<<P001072>%
Ошибка времени компиляции для указания <<<P001295>{} пункт
для класса <<P000965>.

<<P000776>
<суперкласс> ::= <<<P001296>< <typeNotVoid> <mixins>
<<P001297> <mixins>

<mixins> ::= \WITH{{<typeNotVoidList>>
<<P000794>

<<<P001073>%
Область применения \EXTENDS{} и \WITH{} класс <<<P000233> это
Типовой параметр <<<P000234>.

<<<P001074>%
Ошибка времени компиляции, если тип
в <<<P001301>{} класс <<<P000235> это
переменная типа (<<P000827>),
псевдоним типа, не обозначающий класс (<<P000828>),
Перечисленный тип (<<P000829>),Отложенный тип (<<P000830>), тип \DYNAMIC{} (<<P000831>),
или типа <<<P000966> <<<P000237> (P000832).

<<<P001303>{%>
Обратите внимание, что <<<P001304>{} является зарезервированным словом,
что подразумевает, что те же ограничения применяются для типа <<<P001305>,
Аналогичные ограничения установлены для других типов, таких как:
<<P000967> (<<P000833>) и
<<P000968> (<<P000834>)%
?

<<<P001306>{%>
Параметры типового класса доступны в
лексический охват статьи о суперклассе,
Потенциально затеняющие классы в окружающем пространстве.
Поэтому следующий кодекс является незаконным.
и должно вызывать ошибку времени компиляции: %
?

<<P000000>

<<<P001075>%
% % TODO(Eernst): Подумайте о замене всех случаев «суперкласса».
«прямой или косвенный суперкласс», потому что это слишком запутанно.
Класс <<<P000238> <<<P000917> класса <<<P000239> Если {f} либо:

<<P000777>
<<P001307> <<P000240> является суперклассом <<<P000241>, или
<<P001308> <<P000242> является суперклассом класса <<<P000243>,
<<P000244> Это суперкласс P000245.
<<P000795>

<<<P001076>%
Ошибка времени компиляции в классе <<<P000246> Это суперкласс сам по себе.


<<<P001309> Наследование и преобладание
<<P001161>

<<<P001077>%
<<<P000247> Будьте классом, пусть <<<P000248> быть суперклассом <<<P000249>, и
<<<P000250> быть суперклассами <<<P000251> Это также подклассы <<<P000252>.
<<P000253> <<P000918> все конкретные, доступные экземпляры <<<P000254>
которые не были отменены конкретной декларацией в <<P000255>
или по меньшей мере в одном из <<<P000256>.

<<<P001310><%
Было бы более привлекательно дать чисто местное определение наследственности.
Это зависело только от членов прямого суперкласса <<<P000257>.
Класс <<<P000258> может наследовать член <<<P000259> это
Не является членом своего суперкласса. <<P000260>.
Это может произойти, когда член <<<P000261> является частным для библиотеки <<<P000262> <<<P000263>,
тогда как <<P000264> Из другой библиотеки <<P000265>,
Сеть суперклассов <<P000266> включает класс, заявленный в <<<P000267>.%
?

<<<P001078>%
Класс может переопределить членов экземпляра
В противном случае они были бы унаследованы от суперкласса.

<<<P001079>%
<<<P000268> быть классом, объявленным в библиотеке <<<P000269>, и
<<<P000270> быть набором всех суперклассов <<<P000271>,
где <<P000272> Это суперкласс <<<P000273> для <<<P000274>.
<<<P001311><<<P000275> Это встроенный класс <<P000969>.
<<<P000276> декларировать конкретный член <<<P000277>, и
<<<P000278> быть конкретным членом <<<P000279>, который имеет
Имя: <<P000280>,
<<<P000281>> Доступен по адресу <<<P000282>.
Далее <<P000283> переопределение <<P000284>
Если <<P000285> не перевешивается конкретным членом
По крайней мере, один из <<<P000286>
ни того, ни другого <<<P000287> или <<P000288> являются переменными экземпляров.

<<<P001312>{%>
Инстанционные переменные никогда не превосходят друг друга.
Геттеры и сеттеры, индуцированные переменными экземпляров do.%
?

<<<P001313>{%>
Опять же, было бы предпочтительнее местное определение переопределения.
Но не учитывает конфиденциальность библиотеки.%
?

<<<P001314>{%>
Является ли оверрайд законным или нет, определяется относительно
все прямые суперинтерфейсы, а не только интерфейс суперкласса;
Это описано в другом месте
(P000835).
Статические члены никогда ничего не отменяют,
Они могут участвовать в некоторых конфликтах.
Включение деклараций в суперинтерфейсы
(<<P000836>)%
?

<<<P001315>{%>
Для удобства, вот краткое изложение соответствующих правил,
Использование слова «ошибка» для обозначения ошибок времени компиляции.
Помните, что это не нормативно.
Язык управления находится в соответствующих разделах спецификации.

<<P000778><<<P001316> Существует только одно пространство имен
для геттеров, сеттеров, методов и конструкторов (<<P000837>).
Нелокальная переменная <<<P000289> Вводит геттер <<P000290>,
и нелокальной переменной <<<P000291>
Также вводится сеттер
Если оно не является постоянным и окончательным,
или оно является поздним и окончательным и не имеет начального выражения
<<P000970> (<<P000838>, <<P000839>.
Когда мы говорим об участниках, мы имеем в виду
доступные экземпляры, статические или библиотечные переменные,
Геттеры, сеттеры и методы
(P000840).
<<P001317> Вы не можете иметь двух членов с одинаковым именем в одном классе.
Они объявили или унаследовали (<<P000841>). <<P000842>.
<<P001318> Статические члены никогда не наследуются.
<<P001319> Это ошибка, если у вас есть статический элемент под названием <<<P000293> В вашем классе
и экземпляр с тем же именем
(<<P000843>).
<<P001320> Это ошибка, если у вас есть статический установщик <<<P000971>,
и экземпляр член <<<P000295> (P000844).
<<P001321> Это ошибка, если у вас есть статический геттер <<<P000296>
и набором экземпляров <<<P000972> (P000845).
<<P001322> Если вы определили экземпляр члена по имени <<<P000298>,
и ваш суперкласс имеет одноимённый экземпляр,
Они преобладают друг над другом.
Это может быть или не быть законным.
<<P001323> <<P000914>
Если два члена перевешивают друг друга,
Это ошибка, если только она не является правильной.
(P000846).
<<P001324> Сеттеры, геттеры и операторы никогда не имели
факультативные параметры любого вида;
Это ошибка (<<P000847>, <<P000848>, <<P000849>.
<<P001325>
Это ошибка, если член имеет то же имя, что и его класс.
(P000850).
<<P001326> Класс имеет неявный интерфейс (<<P000851>).
<<P001327> Члены суперинтерфейса не наследуются классом.
Они наследуются по неявному интерфейсу.
Интерфейсы имеют собственные правила наследования
(P000852).
<<P001328> Член является абстрактным, если
не имеет тела и не помечена <<P001329>
(<<P000853>, <<P000854>.
<<P001330> Класс является абстрактным, если он явно обозначен <<<P001331>.
<<P001332> Это ошибка, если конкретный класс не реализует какой-либо член.
его интерфейса, и нет нетривиального <<P000973>
(P000855).
<<P001333> Ошибочно называть нефабрику конструктором абстрактного класса.
использование выражения создания экземпляра (<<<P000856>),
Такой конструктор может быть вызван только от другого конструктора.
Использование суперинвокации (<<P000857>).
<<P001334> Если класс определяет экземпляр члена по имени <<<P000299>,
и любой из его суперинтерфейсов имеет подпись участника под названием <<<P000300>,
Интерфейс класса содержит <<<P000301> от самого класса.
<<P001335> Интерфейс наследует всех членов своего суперинтерфейса.
Они не перегружены и не являются членами нескольких суперинтерфейсов.
<<P001336> Если несколько суперинтерфейсов интерфейса
Определение члена с таким же названием, как <<<P000302>,
В большинстве случаев один из членов наследуется.
Этот член (если он существует) является тем, чей тип является подтипом.
всех остальных.
Если такого члена нет, возникает ошибка.
(P000858).
<<P001337> Правило <<<P000859> относится к интерфейсам, а также классам
(P000860).
<<P001338> Это ошибка, если конкретный класс не имеет реализации.
Для метода в его интерфейсе
если он не имеет нетривиального <<P000974>
(P000861).
<<P001339> Идентификатор поименованного конструктора не может быть таким же, как
Базовое имя статического члена, объявленного в том же классе
(<<P000862>).
<<P000796>%
?


<<<P001340>> Суперинтерфейсы
<<P001162>

<<<P001080>%
Класс имеет набор <<<P000919>.Этот набор содержит интерфейс своего суперкласса.
и интерфейсы классов, указанных в
<<<P001341> Положение о классе.

<<P000779>
<интерфейсы> ::= \IMPLEMENTS{{<typeNotVoidList>>
<<P000797>

<<<P001081>%
Область применения <<<P001343> класс <<<P000303> это
Типовой параметр <<<P000304>.

<<<P001082>%
Ошибка времени компиляции, если элемент
в списке типов <<<P001344> класс <<<P000305> это
переменная типа (<<<P000863>),
псевдоним типа, не обозначающий класс (<<<P000864>),
Перечисленный тип (<<P000865>),
Отложенный тип (<<P000866>), тип <<<P001345>< (<<P000867>),
или типа <<<P000975> <<<P000307> (P000868).
Ошибка времени компиляции, если два элемента в списке типов
% % TODO(Eernst): См. понятие nnbd «один и тот же тип».
<<<P001346>> класс <<<P000308> обозначает один и тот же тип <<P000309>.
Это ошибка времени компиляции, если суперкласс класса <<<P000310> это
Один из элементов списка типов <<<P001347> Положение <<P000311>.

<<<P001348>{%>
Можно утверждать, что это безвредно.
повторение типа в списке суперинтерфейсов;
Так зачем делать это ошибкой?
Дело не столько в том, что ситуация ошибочна,
Но это бессмысленно.
Таким образом, это признак того, что программист вполне мог иметь в виду.
Сказать что-то еще - и это ошибка, которая должна быть
Привлекает к себе внимание.%
?

<<<P001083>%
Ошибка времени компиляции в классе <<<P000312> Имеет два суперинтерфейса
Это разные экземпляры одного и того же общего класса.
<<<P001349>{%>
Например, класс не может иметь
оба <<P000976> и <<P000977> в качестве суперинтерфейсов
прямо или косвенно.%
?

<<<P001084>%
При общем классе <<<P000313> объявляет параметр типа <<<P000314>,
Ошибка времени компиляции <<<P000315> происходит в нековариантном положении
% Можно сказать «прямой суперинтерфейс», но легко понять, что это
% достаточно для проверки прямых суперинтерфейсов, и это верно
% для косвенных.
в типе, который определяет суперинтерфейс <<<P000316>.
<<<P001350>{%>
Например, класс не может иметь
\code{Список\VOID{} <<<P001353>(<<P000317>>)
в его \EXTENDS{} или <<P001355>>>{} клаузулы.%
?

<<<P001085>%
Это ошибка времени компиляции, если интерфейс класса <<<P000318> это
Суперинтерфейс самого себя.

<<<P001356>{%>
Класс не наследует членов от своих суперинтерфейсов.
Однако его неявный интерфейс делает.%
?


<<P001357> Конфликты членов класса.
<<P001163>

<<<P001086>%
Некоторые пары деклараций членов класса, смеси, перечня и расширения
не может сосуществовать,
Хотя они не вводят одно и то же имя в один и тот же объем.
В этом разделе указаны эти ошибки.

<<<P001087>%
<<<P000920> из геттера или метода под названием <<<P000319> означает <<<P000320>;
Базовое имя сеттера <<<P000978> является <<<P000322>.
Базовое имя оператора <<<P000323> - <<<P000324>,
За исключением оператора <<P000979> Имя пользователя: <<P000980>.

<<<P001088>%
<<<P001358> Будьте классом.
Ошибка времени компиляции <<<P000325>
Объявлен конструктор по имени <<P000981> и
Статический элемент с базовым именем <<<P000328>.
Ошибка времени компиляции <<<P000329>
Объявляет статический элемент с именем базы <<P000330> и
Интерфейс <<<P000331> имеет элемент экземпляра с именем базы <<<P000332>.
Это ошибка времени компиляции, если интерфейс <<P000333>
имеет метод экземпляра под названием <<<P000334> Настройщик экземпляров с именем базы <<<P000335>.
Ошибка времени компиляции <<<P000336> Статический метод, названный <<P000337>
и статический сеттер с именем базы <<<P000338>.

<<<P001089>%
Когда <<<P001359>{C} представляет собой смесь или расширение,
Ошибки времени компиляции происходят по тем же правилам.
<<<P001360>{%>
В некоторых случаях это избыточно.Например, это уже ошибка для смесителя объявить конструктор.
Но существуют и полезные случаи, например, конфликт между статичным членом.
и примерный член.%
?

<<<P001090>%
Эти ошибки возникают, когда геттеры или сеттеры определены явно.
а также когда они вызваны переменными декларациями.

<<<P001361>{%>
Обратите внимание, что другие ошибки, которые похожи по своей природе, рассматриваются в другом месте.
Например, если <<P000339> Класс, который имеет два суперинтерфейса <<P000340> и <<P000341>,
где <<P000342> Метод получил название <<<P000343>
<<P000344> имеет геттер с именем <<P000345>,
Тогда это ошибка, потому что вычисление интерфейса <<P000346>
включает расчет объединенной подписи члена
(<<P000869>)
из этого геттера и этого метода,
Это ошибка для объединенной подписи члена.
Включая геттер и негетер.%
?


<<<P001362> Интерфейсы
<<P001164>

<<<P001091>%
Этот раздел вводит понятие интерфейсов.
Прежде всего, мы определяем понятие подписей членов.
Эта концепция необходима в определении интерфейсов.

% Нам нужна отдельная концепция подписей членов,
%, так что мы можем получить чистую обработку того, как вычислить
% интерфейс класса и общий интерфейс набора
% суперинтерфейсов, например, суперинтерфейсов класса или
% общий интерфейс, представленный оговоркой «on» в
% «миксин». Например, мы не хотим указывать каждый раз.
% мы проверяем наличие конфликтов, которые не имеют значения для организма;
% что метаданные не имеют значения, интерфейс класса
% никогда не должно содержать тела какого-либо члена в
% первое место. Также четкое разделение синтаксического
Декларация % и подпись участника обеспечивают четкое определение
% возможность введения преобразований, например, для замены
% некоторые типы параметров другими, которые необходимы для
% спецификация правильного отношения переопределения.

<<<P001092>%
<<<P000921> <<P000347>
может быть получено из декларации члена экземпляра класса <<<P000348>.
Он содержит ту же информацию, что и <<P000349>,
За исключением того, что <<P000350> опустить тело, если оно есть;
содержит тип возврата и типы параметров
даже если они подразумеваются в <<<P000351>;
опускает названия позиционных параметров;
не содержит модификатора <<<P001363> по каждому параметру, если таковой имеется;
Опускает метаданные
(<<P000870>);
и опускает информацию о том, является ли член
<<P001364>, <<P001365>, <<<P001366>, или <<P001367>.
Не имеет значения, будет ли <<P000352> дается как явный синтаксис
или индуцируется имплицитно, например, посредством переменной декларации.
Наконец, если <<P000353> имеет формальные параметры,
Каждый из них имеет модификатор <<<P001368>
(<<P000871>)
если и только если этот параметр является ковариантным
(P000872).

<<<P001093>%
Мы используем синтаксис, аналогичный синтаксису абстрактной декларации участника.
Указать подписи членов.
Разница в том, что названия позиционных параметров опущены.
Этот синтаксис используется только для целей спецификации.

<<<P001369>{%>
Подписи участников являются синтетическими объектами, то есть
Они не поддерживаются как конкретный синтаксис в программе Dart.
Это вычисленные объекты, используемые во время статического анализа.
Однако полезно уметь указать на
Свойства подписи члена в данной спецификации
через синтаксическое представление.
Подпись делает его явным
является ли параметр ковариантным по декларированию,
Но остается неявным, является ли он ковариантным по классу.
(P000873).
Причина этого заключается в том, что правило для определения того,
Данное отношение override правильно
(<<P000874>)
зависит от первого, а не от второго.
?

<<<P001094>%
Пусть <<<P001370>{m} будет сигнатурой метода формы

<<P001371>
<<P000982>

<<P001372>
<<<P001373><<<P001374><<<P001375><<<P001376><<<P001377> <<<P001608> T}{1}{n},}

<<P001378><<<P001379><<<P001380>\qquad <<<P001382><<<P001383> <<<P001609> T {\displaystyle T} d}{n+1}{n+k})}.

<<P001384>
<<<P001385>{функциональный тип}{методическая подпись! тип функции
<<P000355> тогда

<<P001386>
<<<P001387>{T_0}.

<<<P001095>%
Пусть <<<P001388>{m} будет сигнатурой метода формы

<<P001389>
<<P000983>

<<P001390>
<<<P001391><<<P001392><<<P001393><<<P001394><<<P001395> <<<P001610> T}{1}{n},}

<<P001396>
<\code\qquad\qquad\{\TripleList\COVARIANT<<<P001612> T}{x}{= d}{n+1}{n+k}<<< P001613>>>).

<<<P001402>
<<<P001403><<<P000357><<<P000357> тогда

<<<P001404>
<<<P001405>{T_0}.

<<<P001096>%
Пусть \DefineSymbol{m} будет сеттерской подписью формы
<<P000984>.
<<<P001407>{функциональный тип} <<P000360> тогда
\FunctionTypeSimple{\VOID}{$T$}.

<<<P001097>%
Тип функции подписи члена остается неизменным, если
Некоторые или все значения по умолчанию опущены.

<<<<P001410><%
Мы не указываем тип функции подписи Геттера.
Для таких подписей мы вместо этого будем напрямую ссылаться на тип возврата.
?

<<<P001098>%
<<<P000922> Синтетическая сущность, которая определяет
Как можно взаимодействовать с объектом.
Интерфейс имеет сигнатуры метода, Getter и Setter.
и набор суперинтерфейсов,
Это опять же интерфейсы.
Каждый интерфейс является неявным интерфейсом класса.
В этом случае мы называем это
<<<P001411>{class interface}{interface!class}
или сочетание нескольких других интерфейсов,
В этом случае мы называем это
<<<P001412>{комбинированный интерфейс}{интерфейс!комбинированный}.

<<<P001099>%
<<<P000362> Будьте классом.
<<<P000923> <<P000363> <<P000364> Интерфейс, который заявляет
подпись, полученная от
каждый экземпляр, заявленный <<<P000365>.
<<P000924> <<P000366> Это прямые суперинтерфейсы <<P000367>
(P000875).

<<<P001413>{%>
Мы говорим, что интерфейс класса "объявляет" эти подписи членов,
Таким образом, мы можем сказать, что интерфейс «объявляет» или «имеет» участник,
Так же, как и для классов.
Обратите внимание, что подпись участника <<P000368> интерфейса класса <<<P000369>
может иметь параметр <<<P000370> с модификатором <<<P001414>,
Хотя <<<P000371> была получена из декларации <<P000372> в <<<P000373>
и параметр, соответствующий <<P000374> в <<<P000375> не
Есть такой модификатор.
Это связано с тем, что <<P000376> может иметь «наследственное» " "
свойство быть ковариантным по декларации
Один из его суперинтерфейсов
(<<P000876>)%
?

<<<P001100>%
С целью выполнения статических проверок по обычным методам вызовов
(<<P000877>)
и извлечение имущества
(<<P000878>),
любой тип <<P000377> Что такое P000378 ограниченный
(<<P000879>),
где <<P000379> класс с интерфейсом <<<P000380>,
Также считается, что интерфейс <<<P000381>.
Точно так же, когда <<<P000382> <<<P000383> где <<P000384> является функциональным типом,
<<P000385> Считается, что метод называется \CALL{} с подписью <<P000386>,
Тип функции <<<P000387> является <<<P000388>.

<<<P001101>%
<<<P001416>{I, \List{I}{1}{k}}%
<<<P000925> <<P000389> Список интерфейсов <<<P001418 >> I}{1}{k)
является интерфейсом, который объявляет набор подписей членов <<<P000390>,
где <<<P000391> определяется как указано ниже.
<<<P000926> <<<P000392> Набор \List I}{1}{k}.

<<<P001102>%
<<<P000393> быть совокупностью всех подписей членов, объявленных <<<P001420> I}{1}{k}.
<<<P001421>{M} - наименьший набор, удовлетворяющий следующему:

<<P000780>
<<<P001422> Для каждого имени <<<P001015> Библиотека <<P000394> <<<P000395> содержит
подпись участника под названием <<<P001016> который доступен по адресу <<P000396>,<<<P000397> быть объединенной подписью участника под названием <<<P001017>
из \List{I}{1}{k} в отношении <<<P000398>.
Это ошибка времени компиляции
если вычисление этой объединенной подписи не удалось.
В противном случае <<<P000399> Содержит <<P000400>.
<<P000798>

<<<P001424>{%>
Интерфейсы должны содержать недоступные подписи участников.
Они могут быть доступны из интерфейсов, связанных с
Декларации подтипов.%
?

<<<P001425>{%>
Например, класс <<<P000401> В библиотеке <<P000402> может объявлять частного члена по имени
<<P000985>,
Класс <<<P000403> В другой библиотеке <<P000404> может расширяться <<<P000405>,
Класс <<P000406> В библиотеке <<P000407> может расширяться <<<P000408>;
<<P000409> может объявить члена, который переопределяет <<P000986> из <<P000410>,
и это отношение переопределения должно проверяться на основе интерфейса <<<P000411>.
Таким образом, мы не можем разрешить интерфейс <<<P000412>
«Забыть» недоступных членов <<P000987>.

Для конфликтов ситуация еще более требовательна:
Классы <<P000413> и <<P000414> В библиотеке <<P000415> могут объявлять частных
<<P000988> и <<P000989>,
Подтип <<<P000416> В другой библиотеке <<<P000417> может иметь
a <<<<P001426> Перечисление обоих <<P000418> и <<P000419>.
В этом случае мы должны сообщать о конфликте, даже если конфликтующие стороны
Заявления не доступны для <<<P000420>,
потому что эти подписи членов затем не пересылаются
(<<P000880>),
и призыв к <<<P000990> на примере $D$ в $L$
должен вернуть "int" в соответствии с подписью первого члена;
и он должен возвращать объект функции согласно второму,
и призыв к <<<P000991>
Должен вернуть <<<P000992> с первой подписью, и она должна быть
(время компиляции или, для динамического вызова, время выполнения) со вторым.%
?

<<<P001427>{%>
Возможно, невозможно одновременно удовлетворить эти ограничения.
Это неизбежно будет сложная семантика.
Поэтому мы решили сделать это ошибкой.
К сожалению, добавление частной декларации
В одной библиотеке может быть нарушен существующий код в другой библиотеке.
Но следует отметить, что конфликты могут быть выявлены локально.
в библиотеке, где существуют частные декларации,
Они возникают только для частных лиц, имеющих
одно и то же имя и несовместимые подписи.
Переименовать этого частного члена во что-нибудь, что не использовалось в этой библиотеке
Устранит конфликт и не сломает ни одного клиента.
?


<<<P001428> Комбинированные подписи членов.
<<P001165>

<<<P001103>%
Этот раздел определяет, как вычислить подпись участника, которая
должным образом обозначать приоритетный набор из нескольких подписей членов;
взяты из заданного списка интерфейсов.

<<<P001429>{%>
Комбинированная подпись имеет тип, который
Подтип всех типов, указанных для этого члена.
Это необходимо для того, чтобы гарантировать, что тип
Член <<<P001018> класса <<<P000423> является четко определенным,
Даже в том случае, когда <<P000424> наследник
несколько различных деклараций <<<P001019>
и не переопределяет <<<P001430>.
В случае неудачи он служит для конкретизации ситуаций.
разработчик должен добавить декларацию, чтобы разрешить двусмысленность.
Подписи участников являются приоритетными в том смысле, что мы выберем
подпись участника интерфейса с наименьшим возможным индексом;
в случае, когда несколько подписей членов одинаково подходят
быть выбрана в качестве объединенной подписи.
То есть «выигрывает первый интерфейс».%
?

<<<P001104>%
Для целей вычисления объединенной подписи члена,
Нам нужно особое понятие
\IndexCustom{равенство}{равенство подписей членов}
подписей членов.
Две подписи членов <<<P000425> и <<P000426> равнымиЕсли они имеют одно и то же имя,
доступ к одному и тому же набору библиотек;
имеют одинаковый тип возврата (для геттеров),
или один и тот же тип функции и одни и те же случаи <<<P001432>
(для методов и сеттеров).

<<<P001433>{%>
В частности, частные методы из разных библиотек никогда не бывают равными.
Лучшие типы также различаются.
Например,
\FunctionTypeSimple{<<P001435>>>}{} и <<P001436>>>{<<P000993>>>}{}
Они не равны, хотя и являются подтипами друг друга.
Мы нуждаемся в этом различии, потому что управление несоответствиями высшего типа является
одна из целей вычисления комбинированного интерфейса.%
?

<<<P001105>%
<<<P001179>%
Теперь мы определили состав подписей.
Пусть <<<P001020> быть идентификатором <<<P000427> библиотека,
<<<P001437>{I}{1}{k} список интерфейсов,
и <<<P001438>{M_0} набор
Все подписи участников от \List{I}{1}{k}
Доступен по адресу <<<P000428>.
The
<<<P001440>{комбинированная подпись участника
<<<<P001022>> \List{I}{1}{k} в отношении <<P000429>>>}{%
объединенная подпись члена
является подписью участника, которая получается следующим образом:

<<<P001106>%
Если <<<P000430> пуст, вычисление объединенной подписи члена не удалось.

<<<P001107>%
Если <<P000431> содержит ровно одну подпись <<P000432>,
Комбинированная подпись участника <<<P000433>.

<<<P001108>%
В противном случае <<P000434> содержит более одной подписи
\DefineSymbol{\List{m}{1}{q}}

<<<P001109>%
<<<P001444> Неудачные смеси
Если <<P000435> содержит по меньшей мере одну подпись
и, по меньшей мере, одной подписью,
Вычисление объединенной подписи не увенчалось успехом.
<<<P001445>

<<<P001110>%
<<<P001446> Геттерс.
Если <<P000436> содержит только геттерские подписи,
расчет объединенной подписи члена продолжается, как описано ниже
Методы и методы,
За исключением того, что он использует тип возврата подписи геттера.
где методы и установщики используют тип функции подписи участника.
<<P001447>

<<<P001111>%
<<<P001448> Методы и установки
В данном случае <<<P000437> состоит только из подписей,
или только подписи,
Потому что имя <<<P001023> в первом случае всегда заканчивается на <<<P001449>
Что никогда не бывает правдой в последнем случае.

<<<P001112>%
<<<P001180>%
Определить, существует ли непустой набор <<<P000438> такой, что
для каждого <<P000439>,
Тип функции <<<P000440> является подтипом
Тип функции <<<P000441> Для каждого <<P000442>.
%
Если такого набора не существует, вычисление объединенной подписи члена не удается.
<<<P001450>{%>
Полезная интуиция в этой ситуации заключается в том, что
Подписи участников не согласованы
Какой тип подходит для члена по имени <<<P001451>.
В противном случае у нас есть набор подписей членов, которые являются наиболее конкретными. " "
в том смысле, что их типы функций являются подтипами их всех.
?

{% Сфера действия для N_min, M_min, N_firstMin.

\def\Nall{\ensuremath{N_{\mbox{\scriptsize{}all}}}}
\def\Mall{<<P001459>>>{M_{\mbox{\scriptsize{}all}}}}
\def\MallFirst{<<P001464>>>{M_{<<P001465>>>{\scriptsize{}first}}}}

<<<P001113>%
В противном случае, когда набор <<<P000443> Как указано выше, существует
Пусть <<<P001467>{} будет наибольшим набором, удовлетворяющим требованию <<<P000444>,
<<<P001181>%
И пусть <<P000445>.
<<<P001468>{%>
То есть, \Mall{} содержит все подписи участников, названные <<<P001024>
с наиболее специфическим типом.
Dart subtyping - это частичный предварительный заказ.
который обеспечивает существование такого наибольшего набора наименьших элементов,
При наличии непустого набора элементов.
Мы можем иметь несколько таких подписей, потому что подписи членов
Они могут быть такими, что не равны.
И все же их типы функций являются подтипами друг друга.
Мы должны вычислить одну подпись участника из <<<P001470>,и мы делаем это, используя порядок заданных интерфейсов.%
?

<<<P001114>%
<<<P001182>%
<<<P000446> наименьшее число, такое, что
<<P000447> Он не пустой.
<<<P000448> Один элемент, который <<<P001471> содержит.
<<<P001472>{%>
Этот набор содержит ровно один элемент, потому что он не пуст.
и ни один интерфейс не содержит более одной подписи участника под названием <<<P001473>.
Другими словами, мы выбираем <<<P000449> В качестве подписи от
Первый возможный интерфейс
Среди наиболее специфических подписей <<<P001474>>%
?
?

<<<P001115>%
Комбинированная подпись участника <<<P000450>,
который получен из <<<P000451> Добавление модификатора <<<P001475>{}
для каждого параметра <<P000452> (если его еще нет)
когда существует <<P000453>
параметр, соответствующий <<<P000454>
(<<P000881>)
Имеет модификатор <<<P001476>.
<<<P001477>{%>
Другими словами,
каждый параметр в объединенной подписи члена помечается ковариантным
если любой из соответствующих параметров отмечен ковариантным,
И не только среди самых известных,
Но среди подписей \emph{all} с именем \id{} (перенаправлено с «P000455»)
в приведенном списке интерфейсов.%
?
<<P001479>


<<<P001480> Суперинтерфейсы
<<P001166>

<<<P001116>%
Интерфейс имеет набор прямых суперинтерфейсов
(P000882).
Интерфейс <<<P000456> представляет собой <<<P000927> интерфейс <<<P000457>
Если {f} либо <<<P000458> Это прямой суперинтерфейс <<P000459>
<<<P000460> является суперинтерфейсом прямого суперинтерфейса <<<P000461>.

<<<P001117>%
Когда мы говорим, что тип <<<P000462>
<<<P001481> {implements}{type!implements a type}
другой тип <<P000463>,
Это означает, что <<<P000464> является суперинтерфейсом <<<P000465>,
<<<P000466> <<<P000467> ограниченный каким-либо типом <<P000468>
(<<P000883>),
<<P000469> Это суперинтерфейс P000470.
Предположим, что <<P000471> Это сырой тип
(<<P000884>)
чье заявление объявляет <<P000472> Типовые параметры.
Когда мы говорим, что тип <<<P000473>
<<P001482> {implements}{type!implements a raw type}
<<P000474>,
Это означает, что существуют типы \List{U}{1}{}
<<<P000475> реализация <<<P001484><<<P000476><<<<P001485> {1}{s}>}

<<<P001486>{%>
Обратите внимание, что это не то же самое, что быть подтипом.
Например, <<P000994> реализация <<P000995>,
Но он не реализует <<P000996>.
Точно так же <<P000997> реализация <<P000998>.
Также обратите внимание, что когда <<<P000477> реализация <<P000478>
где <<P000479> не является подтипом <<<P000999>,
<<P000480> не может быть подтипом <<<P001000>.%
?

<<<P001118>%
Предположим, что <<P000481> Тип и <<<P000482> Это сырой тип, такой, что <<<P000483> реализация <<P000484>.
Существуют уникальные типы <<<P001487>{U}{1}{}
<<P000485> реализация <<<P001488><<<P000486><<<<P001489> {1}{s}>}
Затем мы говорим, что <<<P001490>{U}{1}{s} являются
\IndexCustom{аргументы фактического типа <<P000487>>> при $G$}{%
Введите аргументы! типа в сыром виде,
Мы говорим: «P000489». Это
\IndexCustom{$j$-фактический аргумент типа <<P000491>>> при $G$}{%
Введите аргументы! тип в сыром виде, <<<P000493>,
Для любого <<P000494>.

<<<P001493>{%>
Например, аргумент типа <<<P001001> при <<<P001002>
<<<P001003>.
Эта концепция особенно полезна, когда
цепочка прямых суперинтерфейсов от <<<P000495> до <<<P000496>
не просто передает все типы аргументов в неизменном виде, например,
с таким заявлением, как
<<<P001494><<<P001495><<<P001614><<<P001615> C<X,<<P001616>>> Y>\,\,\EXTENDS\,\,B<List<Y>\, Y, \,X> \{\}.%
?


<<<P001497> Наследование и преобладание
<<P001167>

<<<P001119>%<<<P001183>%
<<<P000497> быть интерфейсом и <<<P000498> Будьте библиотекой.
Определить <<<P000499> быть набором подписей членов
\DefineSymbol{m)
Таким образом, все следующие положения:

<<P000781>
<<P001499> <<P000500> Доступен для <<<P000501> и
<<P001500> <<P000502> Это прямой суперинтерфейс <<P000503> и либо
<<P000782>
<<<P001501> <<P000504> подпись участника <<<P000505> или
<<<P001502> <<P000506> является членом <<<P000507>.
<<P000799>
<<P001503> <<P000508> Не переопределяется <<<P000509>.
<<P000800>

<<<P001120>%
Кроме того, мы определяем <<<P000510> быть
Набор подписей членов \DefineSymbol{m'')
Таким образом, все следующие положения:

<<P000783>
<<<P001505> <<P000511> Это интерфейс класса <<<P000512>.
<<<P001506> <<P000513> декларирует подпись участника <<<P000514>.
<<P001507> <<P000515> Он имеет то же название, что и P000516.
<<<P001508> <<P000517> Доступен по адресу <<<P000518>.
<<P001509> <<P000519> Это прямой суперинтерфейс <<<P000520> и либо
<<P000784>
<<<P001510> <<P000521> подпись участника <<<P000522> или
<<P001511> <<P000523> Является членом <<<P000524>.
<<P000801>
<<P000802>

<<<P001121>%
<<<P000525> быть интерфейсом класса <<<P000526> Объявлено в библиотеке <<P000527>.
<<P000528> <<P000928> Все участники <<P000529>
<<<P000530> <<P000929> <<P000531>, если <<P000532>.

<<<P001122>%
Все ошибки времени компиляции, связанные с преобладанием членов экземпляра
в разделе <<<P000885> Кроме того, вы можете переключаться между интерфейсами.

<<<P001123>%
Если вышеприведенное правило вызовет несколько подписей членов
с таким же названием <<<P001026> Чтобы потом наследовать
Наследуется только один член, а именно:
Комбинированная подпись участника под названием <<<P001512>,
Прямой суперинтерфейс
% Это четко определено, потому что <<P000533> Это интерфейс класса.
в текстовом порядке, в котором они декларируются,
В отношении <<P000534>
(P000886).
Это ошибка времени компиляции
если вычисление упомянутой объединенной подписи участника не удается.


<<<P001513> Правильный член переопределяется
<<P001168>

<<<P001124>%
<<<P001184>%
<<<P000535> и <<P000536> подписи участников с таким же названием <<<P001514>.
Далее <<P000537> <<<P000930> $m'$
Если удовлетворяются все следующие критерии:

<<P000785>
<<P001515>
<<P000539> и <<P000540> Это оба метода, оба геттера или оба сеттера.
<<<P001516>
Если <<<P000541> и <<<P000542> Оба являются геттерами:
Тип возврата <<<P000543> Должен быть подтип возвратного типа <<<P000544>.
<<P001517>
Если <<P000545> и <<P000546> Оба метода или оба сеттера:
<<<P000547> Тип функции <<<P000548>
За исключением того, что параметром типа является встроенный класс <<<P001004>
Для каждого параметра <<P000549> который является ковариантным по декларации
(P000887).
<<<P000550> Тип функции <<<P000551>.
<<P000552> Это должно быть подтипом <<<P000553>.

<<<P001518>{%>
Требование подтипа гарантирует, что список аргументов формирует
которые допустимы для вызова метода с подписью <<<P000554>
Допускается также вызов метода с подписью <<<P000555>.
Например, <<P000556> может принимать 2 или 3 позиционных аргумента;
<<P000557> может принимать 1, 2, 3 или 4 позиционные аргументы, но не наоборот.
Это встроенное свойство правил подтипа функции.
?
<<<P001519>
% % TODO(Eernst): Приходите, это предупреждение удалено.
Если <<<P000558> и <<<P000559> Оба метода являются
<<P000560> необязательный параметр <<<P000561>,
<<P000562> является параметром <<<P000563> соответствует <<<P000564>,<<P000565> имеет значение по умолчанию <<<P000566> и <<<P000567> имеет значение по умолчанию <<<P000568>,
<<P000569> и <<P000570> должны быть идентичными,
Происходит статическое предупреждение.
<<P000803>

<<<P001520>{%>
Обратите внимание, что параметр, который является ковариантным по объявлению
должен иметь тип, удовлетворяющий еще одному требованию;
относительно соответствующих параметров во всех суперинтерфейсах,
прямые и косвенные
(P000888).
Мы не можем сделать это требование частью понятия правильных переопределений.
Потому что правильное переопределение касается только
Отношение к одному суперинтерфейсу.%
?


<<<P001521>{Миксины}
<<P001169>

<<<P001125>%
Смесь описывает разницу между классом и его суперклассом.
Смесь либо получена из существующей классификации классов.
или вводится посредством смешанной декларации.
Это ошибка времени компиляции для получения смеси из
класс, который объявляет генеративного конструктора,
или из класса, который имеет суперкласс, отличный от P001005.

<<<P001126>%
Применение миксина происходит, когда один или несколько миксинов смешиваются с
Классовая декларация через \WITH{} (<<P000889>).
Применение миксина может использоваться для расширения класса на раздел <<<P000890>;
Альтернативно, класс может быть определен как приложение для смешивания.
Как описано в следующем разделе.


<<<P001523>{Смешанные классы}
<<P001170>

<<P000786>
<mixinApplicationClass> ::= \gnewline{}
<identifier> <typeParameters>? <mixinApplication> "

<mixinApplication> ::= <typeNotVoid> <mixins> <interfaces>
<<P000804>

<<<P001127>%
Ошибка времени компиляции, если элемент
Список типов <<<P001525> Положение о применении смеси является
переменная типа (<<P000891>),
тип функции (<<<P000892>),
псевдоним типа, не обозначающий класс (<<P000893>),
Перечисленный тип (<<P000894>),
Отложенный тип (<<P000895>),
тип \DYNAMIC{} (<<P000896>),
Тип \VOID{} (<<P000897>),
или типа <<<P001006 >> для любого <<<P000572> (P000898).
Если <<P000573> является типом в <<<P001528> пункт, <<<P001529>{смешение}{тип! смешение
<<P000574> является смесью, полученной из <<<P000575>, если <<<P000576> обозначает класс,
или смеси, введенной <<<P000577>, если <<<P000578> обозначает смешанную декларацию.

<<<P001128>%
<<<P000579> быть заявлением класса смешанного применения формы

<<P000002>

<<<P001129>%
Ошибка времени компиляции <<<P000580> Перечисленный тип (<<P000899>).
Это ошибка времени компиляции, если какая-либо из <<<P000581> является перечисленным типом
(<<P000900>).

<<<P001130>%
Эффект <<<P000582> В библиотеке <<P000583> Называется она «P000584». в
объем <<P000585>, связанный с классом (<<P000901>) определяется пунктом
<<<P001530><<<P000586> \WITH{} <<P000587>, <<<P001532>, <<P000588>>
с именем <<P000589>, как описано ниже.
Если <<P000590> класс также реализует <<<P000591>, <<<P001533>, <<P000592>.
Если {}f декларация класса закрепляется встроенным идентификатором <<<P001534>,
Определяемый класс становится абстрактным классом.

<<<P001131>%
Оговорка формы <<<P001535><<<P000593> <<<P001536>{} <<P000594>, <<<P001537>, <<P000595>
с именем <<P000596> определяет класс следующим образом:

<<<P001132>%
Если существует только один миксин (<<P000597>), то \code <<<P001539>{} <<P000599>
определяет класс, полученный при применении смеси (<<<P000902>)
Смешивание <<<P000600> (<<<P000903>) к классу, обозначаемому
<<P000601> с именем <<P000602>.

<<<P001133>%
Если существует более одного миксина (<<P000603>), то
<<<P000604> класс, определяемый по формуле <<<P001540>{<<<P000605> <<<P001541>< <<P000606>, <<<P001542>, <<P000607>>
с именем <<<P000608>, где <<<P000609> Это новое имя, и сделать <<<P000610> Абстрактно.<<<P001543><<<P000611> \WITH{} <<P000612>, <<<P001545>, <<P000613> Определяет класс, полученный
посредством применения смеси смеси <<<P000614> к классу <<<P000615> с именем <<P000616>.

<<<P001134>%
В любом случае, пусть <<<P000617> быть классной декларацией с
те же конструкторы, суперклассы, интерфейсы и экземпляры, что и
Определенный класс.
Это ошибка времени компиляции, если объявление <<P000618> вызвать
Ошибка времени компиляции.
% TODO (Eernst): Не совсем!
% Мы не хотим супер-вызовов на ковариантных реализациях.
% - ошибки времени компиляции.

<<<P001546>{%>
Это ошибка, например, если <<<P000619> содержит декларацию члена <<P000620>
который перекрывает подпись участника <<<P000621> в интерфейсе <<<P000622>,
но который не является правильным переопределением <<<P000623>
(<<P000904>)%
?

<<P001547> Смешанная декларация Срок
<<P001171>

<<<P001135>%
Смесь определяет ноль или более
<<<P001548>{mixin Member Declaration}{mixin!member Declaration},
ноль или больше
<<<P001549>{необходимый суперинтерфейс}{mixin!required superinterface},
один
<<<P001550>{комбинированный} суперинтерфейс {mixin!combined superinterface}
ноль или больше
<<<<P001551>{реализованные интерфейсы}{mixin!implemented interface}.

<<<P001136>%
Смесь, полученная из классовой декларации:

<<P000003>

<<P001552>
<<<P001007> Требуемый суперинтерфейс
Комбинированный суперинтерфейс,
<<P000624>, <<<P001553>> <<<P000625> в виде реализованных интерфейсов
и, например, члены <<<P001554>{члены} в виде смешанных деклараций.
Если <<P000626> Это дженерик, как и смесь.

<<<P001137>%
Декларация о смешивании вводит смешивание и обеспечивает объем
Статические декларации членов.

<<P000787>
<mixinDeclaration> ::= \MIXIN{} <typeIdentifier> <typeParameters>
<<<P001556>< (<<<P001557>{} <typeNotVoidList>)? <Интерфейсы>.
\gnewline{}\{" (<metadata> <classMemberDeclaration>)* <<<P001626> "
<<P000805>

<<<P001138>%
Это ошибка времени компиляции, чтобы объявить конструктора в смешанной декларации.

<<<P001139>%
Декларация о смешивании без <<<P001008> пункт является эквивалентным
в соответствии с пунктом \code{\ON{} Объект.

<<<P001140>%
Пусть <<<P000627> будет \MIXIN Декларация формы

<<P000004>

<<<P001141>%
Ошибка времени компиляции, если какой-либо из типов <<P000628> <<<P000629>
или <<<P000630><<<P000631> это
переменная типа (<<P000905>),
тип функции (<<<P000906>),
псевдоним типа, не обозначающий класс (<<<P000907>),
Перечисленный тип (<<P000908>),
отложенный тип (<<<P000909>),
тип <<<P001562>< (<<P000910>),
тип <<<P001563>< (<<P000911>),
или типа <<<P001009> для любого <<<P000633> (<<P000912>).

<<<P001142>%
<<<P000634> быть интерфейсом, заявленным в декларации класса

<<P000005>

<<P001564>
где <<<P000635> - новое имя.
Это ошибка времени компиляции для декларации смешивания, если <<<P000636>
Классовая декларация вызовет ошибку компиляции,
<<<P001565>{%>
То есть, если какой-либо член объявлен более чем одним объявленным суперинтерфейсом,
И нет никакой конкретной подписи для этого члена среди супер.
интерфейсы %
}
Интерфейс <<<P000637> называется
<<<P000931> из заявления о смешивании <<<P000638>.
<<<P001566>{%>
Если заявление о смешивании <<<P000639> имеет только один объявленный суперинтерфейс, <<<P000640>
интерфейс супервызова <<<P000641> имеют точно такие же члены
как интерфейс <<<P000642>.%
?

<<<P001143>%
<<<P000643> Интерфейс, который будет определяться декларацией класса

<<P000006>

<<P001567>
где <<P000644> являются заявлениями членов
Декларация о смешивании <<P000645> За исключением того, что все супервызовы лечатся
как если бы <<<P001568>{} было действительным выражением со статическим типом <<P000646>.Ошибка времени компиляции для микширования <<<P000647> Если это <<P000648> класс
Объявление вызовет ошибку компиляции времени,
<<<P001569>{%>
То есть, если необходимы суперинтерфейсы, реализованные интерфейсы
и декларации не определяют согласованный интерфейс,
если любое заявление члена содержит ошибку компиляции
кроме сверхпризыва,
или если супер-вызов не действителен против интерфейса <<<P000649>.%
?
Интерфейс, представленный декларацией смешивания <<<P000650> имеет тот же член
Подписи и суперинтерфейсы <<<P000651>.

<<<P001144>%
Декларация о смешивании <<<P000652> Добавить микс
<<<P001570>> требуется суперинтерфейс <<<P000653>, <<<P001571>, <<P000654>,
<<<P001572>{комбинированный суперинтерфейс} <<P000655>,
\NoIndex{реализованный интерфейс}<<P000656>>>, <<<P001574>, <<P000657>
и члены экземпляра, объявленные в <<<P000658> как <<<P000932>.


<<<P001575>{Приложение миксина}
<<P001172>

<<<P001145>%
Для суперкласса может быть применена смесь, дающая новый класс.

<<<P001146>%
<<<P000659> Будьте классом,
<<P000660> Будьте смешаны с <<<P001576> требуется суперинтерфейс <<<P000661>, <<<P001577>, <<P000662>,
<<<P001578>{комбинированный суперинтерфейс} <<P000663>,
<<<P001579>{реализованные интерфейсы} <<P000664>, <<<P001580>, <<<P000665>>
\metavar{members} как \NoIndex{mixin members declarations},
Пусть P000666 будет именем.

<<<P001147>%
Ошибка времени компиляции: <<<P000667> <<<P000668> <<<P000669> не осуществлять,
прямо или косвенно, все <<<P000670>, <<<P001583>, <<P000671>.
Это ошибка времени компиляции, если любой из <<<P001584>
супер-призыв члена <<<P000672>
<<<P001585>{%>
(например, <<<P001010>, <<<P001011>,
или <<<P001012>, %
?
<<P000673> не имеет конкретного
Использование <<P000674> который является действительным переопределением члена <<<P000675> в
интерфейс <<P000676>.
<<<P001586>{%>
Мы рассматриваем супер-призывы в миксах как
вызовы интерфейса на комбинированном суперинтерфейсе,
Поэтому мы требуем, чтобы суперкласс приложения для смешивания имел
действительные реализации этих членов интерфейса
На самом деле это супер-призыв.%
?

<<<P001148>%
Применение микширования <<<P000677> до $S$ с именем <<P000679> вводит новый
класс, <<<P000680>, с именем <<<P000681>, суперкласс <<<P000682>,
реализованный интерфейс <<<P000683>
и <<<P001587>{members} в качестве членов экземпляра.
Класс <<<P000684> не имеет статических членов.
Если <<P000685> объявляется любой генеративный конструктор, затем приложение
Внедрение генеративных конструкторов на <<<P000686> следующим образом:

<<<P001149>%
Пусть <<<P000687> будет библиотекой, содержащей приложение микширования.
<<<P001588>{%>
То есть библиотека, содержащая пункт \code{$S$ <<<P001590>< <<P000689>>
или пункт <<<P001591><<<P000690> <<<<P001592>< <<P000691>, \ldots, \ $M_k$, <<P000693 >>
Появление приложения микширования.%
?

<<<P000694> Называется он «P000695».

Для каждого генеративного конструктора формы
<<<P001013>
<<<P000702> Это доступно для <<<P000703>, <<P000704> иметь
неявно заявленный конструктор формы

<<P000007>

<<P001594>
где <<P000705> Получено из <<<P000706> заменяя события <<<P000707>,
который обозначает суперкласс по <<<P000708> и <<<P000709> Получено из <<<P000710> через
Замена происшествий <<<P000711> который обозначает суперкласс по <<<P001595>.
Если <<P000712> Конструктор генеративных констов и <<P000713> не объявляют ни одного
переменных экземпляров, P000714 также является конструктором.

<<<P001150>%
Для каждого генеративного конструктора формы
<<<P001014>
<<<P000726>> Это доступно для <<<P000727>, <<P000728> иметьнеявно заявленный конструктор формы

<<<p000008>

<<P001596>
где <<P000729> Получено из <<<P000730> заменяя события <<<P000731>,
который обозначает суперкласс, по <<<P000732>,
<<P000733> Получено из <<<P000734> заменяя события <<P000735>
который обозначает суперкласс по <<<P001597>,
<<P000736>, <<P000737> - постоянная оценка выражения
То же значение, что и P000738.
Если <<P000739> Конструктор генеративных констов и <<P000740> не объявлять ни одного
переменных экземпляров, P000741 также является конструктором.

<<<P001151>%
Для каждого генеративного конструктора формы
<<<P001598><<<P000742><<<P000743><<<P000744> <<<P001599>, <<P000745> <<P000746>, \{$T_{k+1}$$a_{k+1}$=$d_1$ \ldots, <<P000750>>> $a_{k+n}$ = $d_n$\}]
<<P000753> Это доступно для <<<P000754>, <<P000755> иметь
неявно заявленный конструктор формы

<<P000009>

<<<P001601>
где <<P000756> Получено из <<<P000757> заменяя события <<P000758>
который обозначает суперкласс по <<<P000759>,
<<P000760> Получено из <<<P000761> заменяя события <<<P000762>
который обозначает суперкласс по <<<P001602>,
<<P000763>, <<<P000764>, является постоянным выражением, оценивающим
То же значение, что и <<P000765>.
Если <<P000766> Конструктор генеративных констов и <<P000767> не объявляют ни одного
поля, P000768 также является конструктором.

<<<P001152>%
Если <<P000769> не декларирует никаких генеративных конструкторов, кроме приложения
не вызывает неявного стимулирования каких-либо экспедиторских конструкторов.
В частности, он не вызывает по умолчанию конструктора
(<<P000913>).

<<<P001603>{%>
Конструктор по умолчанию будет иметь суперинициализатор
Это попытка вызвать конструктора, которого не существует.
Разработчики не смогут устранить эту ошибку.
?


<<<P001604>> Расширения.
<<P001173>

<<<P001153>%
В этом разделе указаны расширения.
Этот механизм поддерживает декларирование функций.
которые аналогичны используемым методам,
аналогично неметодическим функциям в этом
Объявляются вне целевого класса,
Они решаются статически.
Резолюция основана на том, имеет ли соответствующее продление сферу действия.
и удовлетворяет ли призыв ряду других требований.

